\documentclass[preprint,12pt]{elsarticle}

\usepackage{amsmath}
\usepackage{booktabs}
\usepackage{graphicx}

\journal{ISA Transactions}

\begin{document}

\begin{frontmatter}

\title{Actuation-Aware Sensor Validation and Rectification for Fault-Resilient Closed-Loop Measurement and Control}

\begin{abstract}
Closed-loop automation can convert transient sensor faults into unsafe or unnecessary actuation. This paper proposes AASVR, an actuation-aware sensor validation and rectification method for multivariate measurement/control systems. AASVR combines robust rate/range plausibility gates, state-machine persistence logic, actuator-aware confirmation windows, and bounded rectification before control authorization. The method is evaluated on hydroponic deployment data and benchmarked against conventional filtering and process-monitoring baselines using external closed-loop process-control datasets. Across real and injected faults, AASVR is designed to reduce anomalous pass-through and false actuation while preserving low detection delay and real-time feasibility. The results evaluate whether actuation-aware validation provides a safer measurement-to-control layer for low-cost closed-loop systems.
\end{abstract}

\begin{keyword}
sensor validation \sep data rectification \sep fault detection \sep process control \sep hydroponics \sep actuation
\end{keyword}

\end{frontmatter}

\section{Introduction}
Closed-loop measurement and control systems continuously convert sensor observations into actuator decisions. In low-cost deployed systems, transient sensor faults can therefore become physical interventions rather than only data-quality problems. This article frames hydroponic controlled-environment agriculture as the primary deployment testbed for a general actuation-aware validation method.

\section{Related Work}
\subsection{Sensor faults in closed-loop measurement and control}
\subsection{Filtering and data-rectification methods}
\subsection{Process-control benchmark datasets}
\subsection{Hydroponic CEA as a validation testbed}

\section{Problem Formulation}
\subsection{Measurement stream and trusted state}
\subsection{Sensor anomaly and actuation-risk definitions}
\subsection{Control-aware objective function}

\section{Proposed Method: AASVR}
\subsection{Overview}
\subsection{Robust plausibility gating}
\subsection{State-machine persistence logic}
\subsection{Rectification policy}
\subsection{Actuation authorization}
\subsection{Parameter optimization}
\subsection{Computational complexity}

\section{Materials and Evaluation Protocol}
\subsection{Hydroponic deployment datasets}
\subsection{External benchmark datasets}
\subsection{Synthetic fault injection}
\subsection{Baseline methods}
\subsection{Evaluation metrics}
\subsection{Statistical analysis}

\section{Results}
\subsection{Hydroponic failure-mode characterization}
\subsection{Hydroponic replay benchmark}
\subsection{External benchmark performance}
\subsection{Synthetic fault-injection robustness}
\subsection{Ablation and sensitivity analysis}
\subsection{Runtime feasibility}
\subsection{Agronomic application demonstration}

\section{Discussion}
\subsection{Why actuation awareness matters}
\subsection{Cross-domain generalization}
\subsection{Practical deployment recommendations}
\subsection{Limitations}

\section{Conclusion}
AASVR provides a lightweight actuation-aware measurement validation layer for closed-loop systems. The evaluation is designed to test whether robust measurement validation, state-machine persistence, rectification, and actuation authorization reduce anomaly pass-through and false actuation compared with conventional filtering baselines.

\section*{Data Availability}
An anonymized supplementary archive will provide code, configuration files, processed-data descriptors, and reproducibility scripts. Restricted third-party datasets are not redistributed; scripts are provided for use after researchers obtain access from the original providers.

\section*{Declaration of Competing Interest}
To be completed in the separate submission package.

\section*{Declaration of Generative AI and AI-Assisted Technologies}
To be completed if AI-assisted tools are used in manuscript preparation.

\bibliographystyle{elsarticle-num}
\bibliography{references}

\end{document}

